\begin{table}[htbp]
\centering
\begin{threeparttable}
\caption{CHFS 2011: Risk Attitude Placebo Test}
\label{tab:chfs_risk_placebo}
{\small
\begin{tabular}{lc}
\toprule
 & (1) \\
 & Infl.\ exp. \\
\midrule
\multicolumn{2}{l}{\textit{Panel A: Placebo Outcome}} \\
\addlinespace
Risk attitude & $-0.010$ \\
 & $(0.008)$ \\
 & $[0.211]$ \\
\addlinespace
Observations & 8,173 \\
Province FE & Yes \\
Controls & Yes \\
\bottomrule
\end{tabular}
}
\begin{tablenotes}[flushleft]
\footnotesize
\item \textit{Notes:} The dependent variable is inflation expectation (1--5, higher = more inflationary). Risk attitude should not predict expectations if the diagnostic channel operates through salience rather than risk preferences. Controls include log income and an urban indicator. Heteroskedasticity-robust standard errors are in parentheses; p-values are in brackets. $^{*}p<0.10$; $^{**}p<0.05$; $^{***}p<0.01$.
\end{tablenotes}
\end{threeparttable}
\end{table}
